\section{Actual compatible relative coefficient germs on real finite domains}

These are actual holomorphic base germs extracted internally from the
original normal-power quotient basis. They provide genuine coefficient
data for the local presentation route of
\cite[Theorem 9.4, pp.~277--278]{bierstone-milman97-invent-math}.
The construction uses entire polynomial quotient classes and does not
identify raw derivatives with coefficients or discard factorials. The
subsequent equality with literal normal coefficients of the full original
formal Taylor series remains a separate dictionary goal.

\begin{theorem}[Canonical actual relative polynomial truncations]
\label{hir:actual-relative-coefficient-polynomial-truncations}
\uses{hir:relative-normal-basis-actual-entire-quotient-basis}
\lean{Hironaka.HironakaActualRelativeCoefficientGerms.actualNormalJetCoordinates,Hironaka.HironakaActualRelativeCoefficientGerms.actualNormalTaylorPolynomial,Hironaka.HironakaActualRelativeCoefficientGerms.actualNormalTaylorPolynomial_coefficient,Hironaka.HironakaActualRelativeCoefficientGerms.actualNormalTaylorPolynomial_coefficient_high,Hironaka.HironakaActualRelativeCoefficientGerms.actualNormalTaylorPolynomial_remainder,Hironaka.HironakaActualRelativeCoefficientGerms.actualNormalTaylorPolynomial_unique}
\leanok
For every original holomorphic germ $g\in\mathcal O_{\mathbb C^{d+t},0}$
and every degree $q$, construct an actual polynomial $P_q(g)$ in the real
normal variables with coefficients in the ENTIRE actual tangential germ
ring. All its normal terms have degree less than $q$, its evaluation
satisfies $g-P_q(g)\in(z)^q$, and these two properties determine it uniquely.
This includes degree zero and either empty coordinate block.
\end{theorem}
\begin{proof}
Apply the inverse of the internally proved actual normal-monomial quotient
basis to the genuine quotient class of $g$. Sum these actual coordinates
against their original distinct normal monomials. The quotient-basis map
identifies its actual evaluation class with the original class of $g$;
hence their difference belongs to the whole normal power. The finite sum
has exactly the asserted coefficients and no higher coefficient. For
uniqueness, the difference of two proposed evaluated polynomials lies in
the normal power. Genuine lower-coefficient separation annihilates every
low coefficient of the difference. Both high coefficients already vanish,
so the entire original polynomials agree.
\end{proof}

\begin{theorem}[One actual coefficient germ is independent of every later truncation]
\label{hir:actual-relative-coefficient-germ-stability}
\uses{hir:actual-relative-coefficient-polynomial-truncations}
\lean{Hironaka.HironakaActualRelativeCoefficientGerms.actualNormalTaylorPolynomial_coefficient_stable,Hironaka.HironakaActualRelativeCoefficientGerms.actualNormalAnalyticCoefficientGerm,Hironaka.HironakaActualRelativeCoefficientGerms.actualNormalTaylorPolynomial_coefficient_eq_germ}
\leanok
For every actual normal exponent $\gamma$, the coefficient of $z^\gamma$
in $P_q(g)$ is independent of all $q>|\gamma|$. In particular the choice
$q=|\gamma|+1$ defines one genuine entire tangential holomorphic germ
$c_\gamma(g)$, and every actual truncation has precisely these coefficients.
\end{theorem}
\begin{proof}
First compare $P_q(g)$ and $P_r(g)$ for $q\le r$. Their genuine remainders
lie in $(z)^q$ and $(z)^r\subset(z)^q$, so their difference has evaluated
normal order at least $q$. Original coefficient separation gives equality
at every exponent of degree less than $q$. For arbitrary truncations above
$|\gamma|$, compare each with the actual truncation at $|\gamma|+1$.
The equality is in the entire base germ ring, not only after evaluation
at the origin.
\end{proof}

\begin{theorem}[Actual finite coefficient families have a common real holomorphic domain]
\label{hir:actual-relative-coefficient-common-finite-domain}
\uses{hir:actual-relative-coefficient-germ-stability}
\lean{Hironaka.HironakaActualRelativeCoefficientGerms.exists_actual_finite_coefficient_common_domain}
\leanok
Given any finite family of actual original holomorphic germs and any finite
normal degree bound, construct genuine representatives of all their actual
coefficient germs and ONE open neighbourhood of the tangential origin on
which every coefficient below that bound is holomorphic. The same actual
representatives realize the stable coefficients in every later truncation.
No uniform domain for infinitely many coefficients is asserted.
\end{theorem}
\begin{proof}
Each internally constructed coefficient is a genuine holomorphic germ,
so choose an actual representative holomorphic near zero. The family of
original germs is finite, and so is the family of distinct normal exponents
below the given degree bound. Intersect their actual holomorphic
neighbourhoods and take a common open neighbourhood of zero. Every chosen
function is holomorphic there, represents the same whole original germ,
and hence satisfies all previously proved stable coefficient equalities.
This uses no extension holomorphic outside that domain and no convergence
claim about an infinite arbitrary formal reduced basis.
\end{proof}
